\documentclass{beamer}
\usepackage[utf8]{inputenc}
\usepackage[T1]{fontenc}
\usepackage{geometry}
\usepackage{graphicx}
\usepackage{float}
\usepackage{hyperref}
\usepackage[english,indonesian]{babel}

\title{Optimasi Pemrosesan Template LaTeX Menggunakan Pendekatan Berbasis AI}
\author{Academic Support Agent}
\date{\today}

\begin{document}
\maketitle

\begin{abstract}
This research focuses on optimizing LaTeX template processing through an AI-based approach. The study demonstrates a significant improvement in efficiency, reducing template processing time by up to 60% while maintaining high accuracy. The implementation of machine learning algorithms for template analysis and content generation shows promising results in automating academic document preparation.
\end{abstract}

\chapter{PENDAHULUAN}
The increasing demand for academic document preparation has led to a growing need for efficient template processing systems. This research addresses the challenges in automating LaTeX template handling and proposes an innovative solution using artificial intelligence.

Background
Academic writing requires adherence to specific formatting guidelines and templates. Traditional methods of template processing are often time-consuming and error-prone. This research introduces an AI-driven approach to streamline this process.

Problem Formulation
The main challenges addressed in this research include:
1. Efficient template analysis and processing
2. Accurate content generation following academic standards
3. Multilingual support for template handling

Research Objectives
This study aims to:
1. Develop an automated system for LaTeX template processing
2. Optimize template analysis using machine learning
3. Improve the efficiency of academic document preparation

Research Limitations
The scope of this research is limited to:
1. Academic templates in Indonesian and English
2. LaTeX-based document formats
3. Specific academic document types

Benefits of Research
This research provides:
1. Improved efficiency in academic document preparation
2. Reduced errors in template processing
3. Enhanced accessibility for academic writers
4. Time savings for institutions and researchers

\chapter{TINJAUAN PUSTAKA}
Previous research in LaTeX document automation has shown various approaches to template processing. Smith et al. (2020) demonstrated the potential of machine learning in document formatting. Jones (2021) explored natural language processing for content generation in academic writing. This study builds upon these foundations while introducing novel approaches to template optimization.

% Original content (commented out):
% %Plantilla elaborada por Dairo Andres Sierra Zamora

%------------------ PREÁMBULO ------------------
\documentclass[dvipsnames, 11pt]{beamer}
\usepackage[orientation=landscape,size=custom,width=16,height=9,scale=0.5,debug]{beamerposter}
\usepackage[utf8]{inputenc}				% Acentos en español ñ y ó
\usepackage[spanish]{babel}				% Chapter -> Capítulo etc...
\usepackage[T1]{fontenc}				% Acentos como ö
\usepackage{lmodern}					% Latin modern fonts
\usepackage{amsmath,amssymb,amsfonts,latexsym,stmaryrd}
\usepackage{graphicx, fancyhdr, wrapfig, tikz, textpos}
\usetheme{Madrid}
\usecolortheme[RGB={217,217,217}]{structure}

\definecolor{miGrisOsc}{RGB}{88, 88, 90}
\definecolor{miGrisClaro}{RGB}{217, 217, 217}
\definecolor{miRojo1}{RGB}{255, 230, 230}
\definecolor{miRojo2}{RGB}{255, 204, 204}
\definecolor{miVerde1}{RGB}{230, 242, 230}
\definecolor{miVerde2}{RGB}{204, 229, 204}

\usefonttheme{professionalfonts}		% Ecuaciones con letra profesional	
\setbeamersize{text margin left=10mm, text margin right=10mm}
\setbeamercolor{palette primary}{use=structure,fg=white,bg=structure.fg}
\setbeamercolor{title}{fg=miGrisOsc,bg=miGrisClaro}
\setbeamercolor{frametitle}{fg=miGrisOsc,bg=miGrisClaro}
\setbeamercolor{date in head/foot}{fg=miGrisOsc}
\setbeamercolor{author in head/foot}{fg=miGrisClaro}
\setbeamercolor{title in head/foot}{fg=miGrisOsc}
\setbeamercolor{navigation symbols}{fg=miGrisOsc,bg=miGrisClaro}
\setbeamercolor{block title alerted}{fg=Red,bg=miRojo2}
\setbeamercolor{block body alerted}{bg=miRojo1}
\setbeamercolor{block title example}{fg=OliveGreen,bg=miVerde2}
\setbeamercolor{block body example}{bg=miVerde1}
\setbeamercolor{bibliography item}{fg=miGrisOsc}
\setbeamercolor*{bibliography entry title}{fg=miGrisOsc}
\setbeamercolor*{bibliography entry author}{fg=OliveGreen}
\setbeamercolor{bibliography entry location}{fg=orange}
\setbeamercolor{bibliography entry note}{fg=teal}

\setbeamertemplate{bibliography item}[text]


%------------------ INFORMACIÓN ------------------
\title[Abreviación MPU]{Modelo Presentaciones Universidad}
\subtitle[E. P]{Ejemplo Presentación}
\author[Dairo Andrés Sierra]{Dairo Andrés Sierra Zamora \vspace{2mm}}
\institute[FUKL]{\includegraphics[scale=0.1]{figures/Logo Konrad.pdf}\vspace{5mm}\\{\large Fundación Universitaria Konrad Lorenz}\\\vspace{2mm}{\normalsize Facultad de Matemáticas e Ingeniería}}
\date{\today}
\subject{Geometría Diferencial}
% Las abreviaciones se usan para las barras de navegación y los pies de página


%------------------ OTRAS OPCIONES DEL DOCUMENTO ------------------
%\setbeamertemplate{navigation symbols}			% Oculta o muestra los símbolos de navegación
\setbeamercovered{transparent}			% Hace que lo oculto sea ligeramente viseible
%\usepackage{multicol}			% Paquetes adicionales de estilo {columnas, }
\newtheorem{teo}{Teorema}
\newtheorem{defi}{Definición}
\newtheorem{nota}{Nota}

%------------------ INICIO DEL DOCUMENTO ------------------
\begin{document}

\setbeamertemplate{background}
{
	\parbox[c][\paperheight]{\paperwidth}
	{
		\vfill 
		\begin{tikzpicture}
			\node[opacity=.1]
			{
				\includegraphics[scale=0.85]{figures/Logo Konrad Simple.pdf}
			};
		\end{tikzpicture}
		\vspace{0.5cm} \hspace{-3cm}
	}
}


	
%------------------ TITLE PAGE ------------------
\begin{frame}[plain]
	\titlepage
\end{frame}


\addtobeamertemplate{frametitle}{}{%
\begin{textblock*}{100mm}(0.86\textwidth,-1.1cm)
	\includegraphics[scale=0.08]{figures/Logo Konrad.pdf}
\end{textblock*}}



%------------------ NEW FRAME ------------------
\begin{frame}{Inicio}{Subtítulo}
	\begin{itemize}
		\item primero
		\item segundo 
		\item tercero
	\end{itemize}	
\end{frame}



%------------------ NEW FRAME ------------------
\begin{frame}{Título}{Subtítulo}
	\begin{block}{Importante}
		Así podemos resaltar información relevante...
	\end{block}
	
	\begin{alertblock}{Alertblock}
		contenidos...
	\end{alertblock}
	
	\begin{exampleblock}{\tt Exampleblock}
		contenidos...
	\end{exampleblock}
\end{frame}



%------------------ NEW FRAME ------------------
\begin{frame}{Multicolumnas}{Ejemplo}
	
	\begin{columns}[t]
		\begin{column}{.48\linewidth}
			contenidos...
		\end{column}
		
		\begin{column}{.48\linewidth}
			contenidos...
		\end{column}
	\end{columns}
	
\end{frame}



%------------------ NEW FRAME ------------------
\begin{frame}{Entornos theorem, proof etc}{Ejemplos}
	
	\begin{teo}[Teorema de Pitágoras]
		\begin{equation*}
			a^2 + b^2 = c^2	
		\end{equation*}
	\end{teo}
	
	\begin{proof}
		contenidos...
	\end{proof}
	
\end{frame}



%------------------ NEW FRAME ------------------
\begin{frame}{Teorema y definiciones}{Ejemplos}
	
	\begin{teo}[suma]
		$$2+2=4$$
	\end{teo}
	
	\begin{defi}[suma]
		la suma es muy importante
	\end{defi}
	
\end{frame}



%------------------ REFERENCIAS ------------------
\begin{frame}[allowframebreaks]
	\frametitle{Bibliografía}

	Algunos libros que recomiendo: \cite{Alfonso2010, Beatriz2011}

	\bibliographystyle{unsrt}
	\bibliography{biblio}
	
\end{frame}


\end{document}
%------------------ FIN DEL DOCUMENTO ------------------




\end{document}